\documentclass[11pt]{article}
\usepackage[margin=1in]{geometry}
\usepackage{amsmath,amssymb,amsthm,mathtools}
\usepackage[hidelinks]{hyperref}
\newtheorem{theorem}{Theorem}
\newtheorem{lemma}[theorem]{Lemma}
\newtheorem{proposition}[theorem]{Proposition}
\newtheorem{corollary}[theorem]{Corollary}
\theoremstyle{definition}
\newtheorem{definition}[theorem]{Definition}
\newcommand{\tr}{\operatorname{tr}}
\newcommand{\SU}{\operatorname{SU}}
\newcommand{\Rea}{\operatorname{Re}}
\newcommand{\Ima}{\operatorname{Im}}
\newcommand{\Var}{\operatorname{Var}}
\newcommand{\Cov}{\operatorname{Cov}}
\newcommand{\dd}{\,\mathrm d}
\title{Exact vacuum-weighted Wilson-loop variation on a finite interacting lattice}
\author{}
\date{8 September 2026}
\begin{document}
\maketitle

\section{Objects, orientations, and both Lie metrics}

Fix an integer $n\geq2$, a finite oriented graph with vertex set $\mathsf V$
and one chosen orientation for every physical edge in a nonempty edge set
$\mathsf E$, and a finite list $\mathsf P$ of oriented closed plaquette words.
Write $N=|\mathsf E|$. Reverse traversal of $e$ uses $U_e^{-1}$, not a new
independent link. The configuration manifold and its probability Haar measure
are
\[
 \mathcal Q=\SU(n)^{\mathsf E},\qquad
 d\lambda(U)=\prod_{e\in\mathsf E}dU_e,\qquad \lambda(\mathcal Q)=1.
\]
For a word $C=((e_1,\varepsilon_1),\ldots,(e_\ell,\varepsilon_\ell))$
with $\varepsilon_j\in\{1,-1\}$, define, in the written order,
\[
 H_C(U)=U_{e_1}^{\varepsilon_1}\cdots U_{e_\ell}^{\varepsilon_\ell},
 \qquad W_C(U)=\tr H_C(U).
\]
A simple edge loop here means a nonempty closed word using each physical edge
at most once. Vertices may be repeated; all statements below needing simplicity
need precisely the absence of repeated physical edges. Ordinary simple
plaquettes and primitive winding loops on a sufficiently large periodic
hypercubic lattice satisfy this definition. If a small periodic lattice makes
one plaquette repeat an edge, use the repeated-occurrence formula below.

Choose an anti-Hermitian traceless basis $T_a$, $1\leq a\leq n^2-1$, with
\[
 -2\tr(T_aT_b)=\delta_{ab}.
\]
The derivatives requested in this calculation are exactly
\[
 X_{e,a}f(U)=\left.\frac d{dt}\right|_{t=0}
 f(\ldots,e^{tT_a}U_e,\ldots),\qquad
 \Delta_e^T=\sum_aX_{e,a}^2.
\]
In this convention $X_{e,a}$ is skew-adjoint in $L^2(\lambda)$. Indeed,
left invariance of Haar gives
\[
 \int f(e^{tT_a}U_e)g(U_e)dU_e
 =\int f(U_e)g(e^{-tT_a}U_e)dU_e;
\]
differentiation proves the assertion,
with the complex conjugate placed on the first argument when appropriate.

There are two retained Lie metrics, related by an explicit dictionary:
\begin{align}
 g_T(A,B)&=-2\tr(AB),&
 c(A,B)&=-\tfrac12\tr(AB)=\tfrac14g_T(A,B),\label{eq:metrics}\\
 c(T_a,T_b)&=\tfrac14\delta_{ab},&
 S_a&=2T_a,& c(S_a,S_b)&=\delta_{ab},\nonumber\\
 X^c_{e,a}&=2X_{e,a},&
 \Delta_e^c&=4\Delta_e^T.\label{eq:metric-laplace}
\end{align}
Thus a Casimir eigenvalue measured with $c$ is four times its $g_T$ value.
The Hamiltonian in this paper is the specified operator
\begin{equation}\label{eq:H}
 H=-\sum_{e\in\mathsf E}\kappa_e\Delta_e^T+V,
 \quad
 V(U)=\sum_{p\in\mathsf P}b_p\bigl(n-\Rea W_p(U)\bigr),
 \quad \kappa_e>0,\ b_p>0.
\end{equation}
In the original cusp metric it is identically
$H=-\sum_e(\kappa_e/4)\Delta_e^c+V$.
Keeping the coefficient $\kappa_e$ in front of $\Delta_e^c$ would instead
multiply the kinetic operator by four and would define a different
Hamiltonian. No such replacement is made here. The inequalities
$|\tr U|\leq n$ for $U\in\SU(n)$ and
\begin{equation}\label{eq:Vbounds}
 0\leq V\leq M,\qquad
 M=2nB,\qquad B=\sum_{p\in\mathsf P}b_p,
\end{equation}
retain every plaquette term. The trace bound follows by summing the $n$
unit-modulus eigenvalues. The constant $2nB$ is a valid bound; no assertion
that all plaquettes attain their separate maxima simultaneously is needed.

The probability Haar convention also has an exact volume dictionary.
On the compact configuration manifold, whose real dimension is
$d=N(n^2-1)$, \eqref{eq:metrics} gives
$d\operatorname{vol}_c=2^{-d}d\operatorname{vol}_{g_T}$.
If $\mathcal V_c=\operatorname{vol}_c(\mathcal Q)$ and
$\mathcal V_T=\operatorname{vol}_{g_T}(\mathcal Q)$, then
\[
 d\lambda=\mathcal V_c^{-1}d\operatorname{vol}_c
         =\mathcal V_T^{-1}d\operatorname{vol}_{g_T},\qquad
 \mathcal V_c=2^{-d}\mathcal V_T.
\]
The map $f\mapsto\sqrt{\mathcal V_c}\,f$ is unitary from
$L^2(d\operatorname{vol}_c)$ to $L^2(d\lambda)$ and intertwines the same
differential Hamiltonian. A ground state normalized in the former
space is therefore taken to $\psi_0=\sqrt{\mathcal V_c}\,\psi_c$, so
$\psi_0^2d\lambda=\psi_c^2d\operatorname{vol}_c$ exactly. These are
volume identities on the compact group product only; they do not
identify a noncompact cusp measure with Haar measure.

Gauge transformations are
$(g\cdot U)_e=g_{s(e)}U_eg_{t(e)}^{-1}$ for
$g\in\SU(n)^{\mathsf V}$. Along a closed word all adjacent factors cancel,
so $H_C$ changes by conjugation at its base vertex and $W_C$ is invariant.
Haar measure, each $\Delta_e^T$, and $V$ are gauge invariant. The physically
restricted finite Hilbert space is therefore the closed invariant subspace
$\mathcal H_{\rm inv}\subset L^2(\mathcal Q,\lambda)$, and $H$ preserves it.

\section{Fierz identities with the stated normalization}

\begin{lemma}\label{lem:Fierz}
For arbitrary complex $n\times n$ matrices $A,B$,
\begin{align}
 \sum_a(T_a)_{ij}(T_a)_{kl}
 &=-\frac12\left(\delta_{il}\delta_{jk}
                  -\frac1n\delta_{ij}\delta_{kl}\right),\label{eq:Fierz}\\
 \sum_a\tr(T_aA)\tr(T_aB)
 &=-\frac12\left(\tr(AB)-\frac1n\tr A\tr B\right),\label{eq:bilinear}\\
 \sum_a\tr(T_aA)\overline{\tr(T_aB)}
 &=\frac12\left(\tr(AB^\dagger)
       -\frac1n\tr A\,\overline{\tr B}\right),\label{eq:sesquilinear}\\
 \sum_aT_a^2&=-C_FI_n,\qquad C_F=\frac{n^2-1}{2n}.\label{eq:fundcasimir}
\end{align}
In particular, for $A\in\SU(n)$,
\begin{equation}\label{eq:singlecontraction}
 \sum_a|\tr(T_aA)|^2=\frac{n^2-|\tr A|^2}{2n}.
\end{equation}
\end{lemma}
\begin{proof}
The matrices $Q_0=I_n/\sqrt n$ and $Q_a=i\sqrt2T_a$ form an orthonormal
basis of the real vector space of Hermitian matrices for the scalar product
$\tr(QR)$. Their complex linear span is all complex matrices. Expanding an
arbitrary matrix $D$ gives
$D=\sum_{\alpha=0}^{n^2-1}Q_\alpha\tr(Q_\alpha D)$.
Taking the coefficient of $D_{lk}$ in entry $ij$ proves
$\sum_\alpha(Q_\alpha)_{ij}(Q_\alpha)_{kl}=\delta_{il}\delta_{jk}$.
Subtract the $Q_0$ term and use $(i\sqrt2)^2=-2$ to obtain
\eqref{eq:Fierz}. Contract it with $A_{ji}B_{lk}$ to get
\eqref{eq:bilinear}. Anti-Hermiticity gives
$\overline{\tr(T_aB)}=-\tr(T_aB^\dagger)$; this proves
\eqref{eq:sesquilinear}. Put $k=j$ in \eqref{eq:Fierz} and sum $j$ to
obtain \eqref{eq:fundcasimir}. Finally $AA^\dagger=I_n$ gives
\eqref{eq:singlecontraction}.
\end{proof}

\section{Every edge orientation and the exact overlap morphism}

Fix one occurrence of physical edge $e$ in $C$ and write the ordered word
as $P U_e^{\varepsilon}Q$, where $P$ and $Q$ are the products before and
after that occurrence. Cyclic trace invariance gives the following
matrices based at the original source $s(e)$:
\begin{equation}\label{eq:based}
 A_{C,e,r}=
 \begin{cases}
 U_eQP,&\varepsilon=+1,\\
 QPU_e^{-1},&\varepsilon=-1.
 \end{cases}
\end{equation}
The occurrence index $r$ is suppressed when there is only one occurrence.
Each $A_{C,e,r}$ lies in $\SU(n)$, has trace $W_C$, and transforms by
conjugation with $g_{s(e)}$ under a gauge transformation. These statements
follow either by multiplying the factors or by following the cyclically
based closed path at $s(e)$.

The derivative of an inverse is
$X_{e,a}U_e^{-1}=-U_e^{-1}T_a$, since
$(e^{tT_a}U_e)^{-1}=U_e^{-1}e^{-tT_a}$. Consequently
\begin{align}
 X_{e,a}\tr(PU_eQ)&=\tr(PT_aU_eQ)=\tr(T_aU_eQP),\nonumber\\
 X_{e,a}\tr(PU_e^{-1}Q)&=-\tr(PU_e^{-1}T_aQ)
 =-\tr(T_aQPU_e^{-1}).\label{eq:orientations}
\end{align}
For an arbitrary closed word the product rule now gives
\begin{equation}\label{eq:repeatedderivative}
 X_{e,a}W_C=\sum_{r:e_r=e}\varepsilon_r\tr(T_aA_{C,e,r}).
\end{equation}
All factors in \eqref{eq:based} are evaluated at the original configuration;
there is no independent-link approximation in this formula.

For two words define the edgewise sesquilinear and bilinear contractions
by
\[
 \Gamma_e(F,G)=\sum_a(X_{e,a}F)\overline{X_{e,a}G},\qquad
 \Gamma_e^{\rm bil}(F,G)=\sum_a(X_{e,a}F)(X_{e,a}G).
\]
Lemma \ref{lem:Fierz} proves, occurrence by occurrence,
\begin{align}
\Gamma_e(W_C,W_D)
 &=\frac12\sum_{r:e_r=e}\sum_{s:d_s=e}\varepsilon_r\eta_s
 \left[\tr(A_{C,e,r}A_{D,e,s}^{\dagger})
       -\frac1nW_C\overline{W_D}\right],\label{eq:overlap-general}\\
\Gamma_e^{\rm bil}(W_C,W_D)
 &=-\frac12\sum_{r:e_r=e}\sum_{s:d_s=e}\varepsilon_r\eta_s
 \left[\tr(A_{C,e,r}A_{D,e,s})-\frac1nW_CW_D\right].
 \label{eq:overlap-bil-general}
\end{align}
These identities construct the exact relation between the two loop
variations: $A_{C,e,r}A_{D,e,s}^{-1}$ is the concatenation of the first
based loop and the inverse of the second, and its trace is gauge invariant.
The nonconjugated contraction concatenates the two based loops directly.
An edge not occurring in both words contributes zero.

For simple edge loops, a common edge contributes exactly
\begin{align}
\Gamma_e(W_C,W_D)
 &=\frac{\varepsilon_{Ce}\varepsilon_{De}}2
   \left[\tr(A_{Ce}A_{De}^{\dagger})
                    -\frac1nW_C\overline{W_D}\right],\label{eq:overlap-simple}\\
\Gamma_e^{\rm bil}(W_C,W_D)
 &=-\frac{\varepsilon_{Ce}\varepsilon_{De}}2
   \left[\tr(A_{Ce}A_{De})-\frac1nW_CW_D\right].\nonumber
\end{align}
In particular, for every used edge of a simple loop,
\begin{equation}\label{eq:simpleenergy}
 \Gamma_e(W_C,W_C)=\frac{n^2-|W_C|^2}{2n},\qquad
 \sum_aX_{e,a}^2W_C=-C_FW_C.
\end{equation}
The second equality follows directly from \eqref{eq:orientations} and
\eqref{eq:fundcasimir}; in negative orientation the two derivative signs
cancel. Both equalities are zero on unused edges. Define
$K_C=\sum_{e\in C}\kappa_e$. Then the full pointwise energy is
\begin{equation}\label{eq:simplefullenergy}
 \Gamma_\kappa(W_C,W_C):=\sum_e\kappa_e\Gamma_e(W_C,W_C)
       =\frac{K_C}{2n}(n^2-|W_C|^2).
\end{equation}

For later real plaquette observables the exact companion identity is
\begin{equation}\label{eq:realenergy}
 \Gamma_\kappa(\Rea W_C,\Rea W_C)
 =\frac{K_C}{4}\left[n-\Rea\tr(H_C^2)
                           -\frac2n(\Ima W_C)^2\right].
\end{equation}
To prove it, use $(\Rea z)^2=(|z|^2+\Rea z^2)/2$ in the derivative
sum, apply both contractions in Lemma \ref{lem:Fierz}, and observe
$\tr(A_{Ce}^2)=\tr(H_C^2)$ because cyclic rotations of invertible factors
are conjugate matrices. Thus the right side is nonnegative despite its
written difference of terms. The same energy holds for $n-\Rea W_C$.
Measured using $X^c=2X$, every $\Gamma_e$ in this section is multiplied
by four; pairing it with kinetic coefficient $\kappa_e/4$ recovers
exactly \eqref{eq:simplefullenergy} and \eqref{eq:realenergy}.

\section{Existence, exact ground-state transform, and orthogonality}

The domain of the nonnegative quadratic form of $H$ is the first Sobolev
space on $\mathcal Q$, with the equivalent norm given by the finitely many
fields $X_{e,a}$. Its form is
\[
 q_H(u,v)=\sum_{e,a}\kappa_e\int\overline{X_{e,a}u}X_{e,a}v\,d\lambda
                   +\int V\overline uv\,d\lambda.
\]
We use the self-adjoint operator associated to this closed form.

\begin{lemma}\label{lem:ground}
$H$ has a unique normalized strictly positive ground eigenfunction
$\psi_0\in C^\infty(\mathcal Q)$, with $H\psi_0=E_0\psi_0$ and
$\int\psi_0^2d\lambda=1$. It is gauge invariant. In particular the bottom
eigenvalue of the gauge-invariant restriction is the same $E_0$.
\end{lemma}
\begin{proof}
The free positive operator $H_0=-\sum_e\kappa_e\Delta_e^T$ has a complete
matrix-coefficient eigenbasis by the compact-group Peter--Weyl theorem
\cite{PeterWeyl}; its finite-multiplicity eigenvalues tend to infinity.
The latter fact also follows from the explicit weights and bounds proved
in Section \ref{sec:heat}. Its resolvent is compact. Bounded real
multiplication by $V$ preserves self-adjointness and compact resolvent:
for sufficiently negative $z$, the identity
$(H-z)^{-1}=(H_0-z)^{-1}[I+V(H_0-z)^{-1}]^{-1}$ proves compactness by a
convergent Neumann series. Hence the bottom of the spectrum is an
eigenvalue. The inequality $|X|u||\leq|Xu|$ almost everywhere (obtain it
first for $\sqrt{|u|^2+\epsilon^2}$ and then let $\epsilon\downarrow0$)
shows that the modulus of a minimizing eigenfunction is also minimizing.
Choose a nonnegative normalized one. The elliptic equation
$H_0\psi_0=(E_0-V)\psi_0$ and local elliptic regularity, iterated with
smooth $V$, give $\psi_0\in C^\infty$.

For completeness, positivity needs no assumption about an unknown vacuum.
The parabolic comparison inequality proved in Section \ref{sec:heat} is
$e^{-tH}f\geq e^{-Mt}e^{-tH_0}f$ for $f\geq0$.
At the explicit time $t_*$ in that section the kernel of $e^{-t_*H_0}$
is at least $2^{-N}$. Thus
$\psi_0(U)\geq e^{(E_0-M)t_*}2^{-N}\int\psi_0d\lambda>0$.
The next theorem, whose calculation uses only this already established
positivity and the eigenvalue equation, proves uniqueness: if $u$ is
another ground eigenfunction, put $f=u/\psi_0$ in its identity. The
result is $\int\psi_0^2\sum_{e,a}\kappa_e|X_{e,a}f|^2d\lambda=0$.
The fields span each tangent space and $\mathcal Q$ is connected, so
$f$ is constant. Gauge transformations preserve $H$, positivity and
normalization, and consequently fix $\psi_0$ by this uniqueness.
\end{proof}

Set
\begin{equation}\label{eq:vacuummeasure}
 d\mu=\rho\,d\lambda,\qquad \rho=\psi_0^2,
 \qquad \mu(F)=\int F\,d\mu.
\end{equation}
This measure is determined by the interacting Hamiltonian
\eqref{eq:H}; no other measure is substituted for it below.
Multiplication $\mathcal U f=\psi_0 f$ is a unitary map from
$L^2(\mu)$ onto $L^2(\lambda)$, since
$\|\mathcal U f\|_\lambda^2=\int|f|^2d\mu$, and division by the positive
$\psi_0$ is its inverse. It also maps the respective gauge-invariant
subspaces onto one another.

\begin{theorem}\label{thm:ground-transform}
For all smooth complex $f,g$,
\begin{align}
 Lg:=\mathcal U^{-1}(H-E_0)\mathcal Ug
 &=-\sum_{e,a}\kappa_e
       \bigl[X_{e,a}^2g+2(X_{e,a}\log\psi_0)X_{e,a}g\bigr]
       \label{eq:groundop}\\
 &=-\rho^{-1}\sum_{e,a}\kappa_e X_{e,a}(\rho X_{e,a}g),\nonumber\\
 \langle\psi_0f,(H-E_0)\psi_0g\rangle_\lambda
 &=\sum_{e,a}\kappa_e\int\overline{X_{e,a}f}X_{e,a}g\,d\mu.
 \label{eq:groundform}
\end{align}
The identities extend to the corresponding closed form domains. For any
nonconstant smooth $F$, the vector
$u_F=\psi_0(F-\mu(F))$ is orthogonal to $\psi_0$, and
\begin{equation}\label{eq:Rayleigh-general}
 \frac{\langle u_F,(H-E_0)u_F\rangle}{\|u_F\|^2}
 =\frac{\mu(\Gamma_\kappa(F,F))}
        {\mu(|F|^2)-|\mu(F)|^2}.
\end{equation}
\end{theorem}
\begin{proof}
The product rule gives
$X^2(\psi_0g)=\psi_0X^2g+2(X\psi_0)Xg+gX^2\psi_0$ for each $X$.
Subtracting $g(H-E_0)\psi_0=0$ from $(H-E_0)(\psi_0g)$ proves
\eqref{eq:groundop}; its divergence form uses $X\rho=2\psi_0X\psi_0$.
Multiply by $\rho\overline f$ and integrate. The skew-adjointness of
$X$ with respect to $\lambda$, proved above, gives
$-\int\overline f X(\rho Xg)d\lambda
=\int\rho\overline{Xf}Xg d\lambda$. Summing proves
\eqref{eq:groundform}. On a compact manifold multiplication and division
by the smooth positive $\psi_0$ are bounded on the first Sobolev space,
so smooth approximation proves the form extension. Finally
$\langle\psi_0,u_F\rangle=\mu(F)-\mu(F)=0$,
$\|u_F\|^2=\mu(|F|^2)-|\mu(F)|^2$, and subtracting a constant leaves
all derivatives unchanged. The variance is strictly positive because
$F$ is nonconstant and $\rho>0$ everywhere. These observations prove
\eqref{eq:Rayleigh-general}.
\end{proof}

In particular, put $m_C=\mu(W_C)$ and $r_C=\mu(|W_C|^2)$. For a simple
edge loop the exact answer is
\begin{equation}\label{eq:Rayleigh-loop}
 Q_C=\frac{K_C}{2n}\,
       \frac{n^2-r_C}{r_C-|m_C|^2}.
\end{equation}
The potential disappears from the displayed Dirichlet integrand because
the exact ground-state equation was used. It remains in $\rho$, $m_C$,
and $r_C$, as well as in the drift term in \eqref{eq:groundop}. Thus
this identity does not set $V$ or any $b_p$ to zero.

\section{Center symmetry, exact selection rules, and character limits}

Assign a central element $z_e\in Z(\SU(n))$ to each physical edge, subject
to
\begin{equation}\label{eq:centercondition}
 \prod_{r\in p}z_{e_r}^{\varepsilon_r}=1
 \quad\hbox{for every retained plaquette }p.
\end{equation}
Define $\sigma_z(U)_e=z_eU_e$. This is a Haar-preserving diffeomorphism,
commutes with the edge derivatives, and preserves every plaquette
holonomy because the factors $z_e$ are central and obey
\eqref{eq:centercondition}. It therefore preserves $H$ and, by uniqueness
and positivity, $\psi_0$ and $\mu$. For every closed word,
\[
 W_C(\sigma_z U)=\zeta_C W_C(U),\qquad
 \zeta_C=\prod_{r\in C}z_{e_r}^{\varepsilon_r}.
\]
Integration proves the exact selection rules
\begin{equation}\label{eq:selection}
 \mu(W_C)=0\quad\hbox{if }\zeta_C\ne1,
 \qquad
 \mu(W_C\overline{W_D})=0\quad\hbox{if }\zeta_C\overline{\zeta_D}\ne1.
\end{equation}
For powers and products the same proof applies with their product charge.

On a periodic hypercubic lattice, multiply the links in one coordinate
direction crossing one fixed periodic cut by $z\in Z(\SU(n))$, and
leave the other links unchanged. Each oriented plaquette crosses that
cut with total signed count zero, which verifies
\eqref{eq:centercondition} directly. A loop with integer winding $q$
in that direction acquires $z^q$. A primitive single winding has a
nontrivial charge for a nonidentity suitable $z$, and therefore $m_C=0$.
A loop whose winding in every direction is divisible by $n$ does not
obtain a zero mean from this argument. Contractible plaquettes are
center neutral; their means must be retained.

For a charged simple loop the quotient becomes
\begin{equation}\label{eq:chargedquotient}
 Q_C=\frac{K_C}{2n}\left(\frac{n^2}{r_C}-1\right),
 \qquad 0<r_C<n^2.
\end{equation}
The strict inequalities follow from $\rho>0$ and the fact that the
holonomy map of a simple loop has image all of $\SU(n)$: fix the other
edges and vary any one used edge. Some images have trace zero and
others have trace $n$, and neighborhoods of both have positive Haar
measure. In particular center symmetry does not set $r_C$ equal to one.

\begin{lemma}\label{lem:Haar-loop}
For every simple edge loop,
\[
 \int W_Cd\lambda=0,\quad
 \int|W_C|^2d\lambda=1,\quad
 \int(n^2-|W_C|^2)d\lambda=n^2-1.
\]
\end{lemma}
\begin{proof}
Condition on every edge except one appearing once. Its loop holonomy
is $A U_e^{\pm1}B$ for fixed group matrices $A,B$, and is Haar distributed
by left and right invariance and invariance under inversion. It suffices
to integrate over one Haar matrix $U$. Its mean trace is zero since
$U\mapsto zU$ preserves Haar for a nonidentity central $z$.
For arbitrary $D$, the matrix $\int UDU^\dagger dU$ commutes with every
$\SU(n)$ matrix by invariance, and so is $(\tr D)I_n/n$.
To see the scalar conclusion without an additional representation
assumption, commuting with diagonal unitary matrices kills the
off-diagonal entries and commuting with determinant-one unitary changes
of coordinate axes makes all diagonal entries equal. Set $D=E_{jl}$
and compare entries to get
$\int U_{ij}\overline{U_{kl}}dU=\delta_{ik}\delta_{jl}/n$.
Taking $j=i$, $l=k$ and summing $i,k$ gives $\int|\tr U|^2dU=1$.
The last asserted integral follows by subtraction.
\end{proof}

There is also an exact character map that retains the second moment.
The conjugation action $A\mapsto UAU^{-1}$ on $\operatorname{End}(\mathbb C^n)$
has trace $\tr U\tr U^{-1}=|\tr U|^2$. The invariant decomposition into
scalar and traceless matrices has characters $1$ and $\chi_{\rm ad}$.
Consequently
\begin{equation}\label{eq:adjointmoment}
 |\tr U|^2=1+\chi_{\rm ad}(U),\quad
 r_C=1+\mu(\chi_{\rm ad}(H_C)),\quad
 -1\leq\chi_{\rm ad}(U)\leq n^2-1.
\end{equation}
This derivation proves the equality and both bounds directly; the adjoint
character is center neutral, so \eqref{eq:selection} imposes no zero
expectation on it.

\begin{proposition}\label{prop:symmetryinsufficient}
Center and gauge invariance, positivity of a smooth density, and the
pointwise character bounds alone imply no strictly positive uniform
lower bound on $r_C$ and no upper bound strictly smaller than $n^2$.
This assertion concerns those premises alone, not the narrower family
of ground states of \eqref{eq:H}.
\end{proposition}
\begin{proof}
Let $h(U)=|W_C(U)|^2$. Its range includes $0$ and $n^2$.
For an explicit trace-zero matrix take
$\zeta=e^{2\pi i/n}$ and
$A=\alpha\operatorname{diag}(1,\zeta,\ldots,\zeta^{n-1})$ with
$\alpha=e^{-\pi i(n-1)/n}$. The geometric sum is zero and
$\det A=\alpha^n\zeta^{n(n-1)/2}=1$. Choose one edge to realize $A$
as holonomy; identity holonomy realizes $n^2$.
For each $s>0$, the densities
\[
 \rho_s^- =\frac{e^{-sh}}{\int e^{-sh}d\lambda},\qquad
 \rho_s^+ =\frac{e^{sh}}{\int e^{sh}d\lambda}
\]
are smooth, strictly positive, gauge invariant, and invariant under
every center transformation of the stated type. For $\epsilon>0$, the
set $h<\epsilon/2$ is nonempty open and has Haar measure $a_\epsilon>0$.
Its contribution to the first denominator is at least
$a_\epsilon e^{-s\epsilon/2}$, whereas the numerator mass of $h\geq\epsilon$
is at most $e^{-s\epsilon}$. Thus this latter probability tends to zero
and $\int h\rho_s^-d\lambda\to0$ by boundedness. Apply the same argument
to $n^2-h$ for $\rho_s^+$ to obtain $\int h\rho_s^+d\lambda\to n^2$.
For a nontrivially charged loop, both measures also have zero loop mean
by the selection-rule proof. No claim is made that these constructed
densities arise from positive Wilson plaquette coefficients in
\eqref{eq:H}; their role is to prove exactly the asserted insufficiency
of symmetry and character bounds alone.
\end{proof}

\section{Explicit all-coupling finite-volume estimates}\label{sec:heat}

The full Hamiltonian supplies more information than those symmetry
premises. The following bounds use it and hold for arbitrary positive
values of all retained couplings. Define
\begin{equation}\label{eq:constants}
 \kappa_* =\min_{e\in\mathsf E}\kappa_e,
 \quad D=n^2-2,
 \quad t_* =\frac{2(D\log2+\log3)}{(n-1)\kappa_*},
 \quad a_* =3^{-2N}e^{-2Mt_*}=3^{-2N}e^{-4nBt_*}.
\end{equation}
These constants are dimensionally consistent: $t_*$ is inverse energy,
$M$ is energy, and $a_*$ is dimensionless. They are estimates, not a
claim of optimal constants.

\begin{lemma}\label{lem:heat}
The integral kernel of $P_t=e^{-tH_0}$ satisfies
\begin{equation}\label{eq:kernelbounds}
 2^{-N}\leq p_{t_*}(U,V)\leq(3/2)^N
 \qquad\hbox{for all }U,V\in\mathcal Q.
\end{equation}
\end{lemma}
\begin{proof}
We give the representation constants to fix the heat time and Lie
metric exactly. The irreducibles of $\SU(n)$ are indexed by
$m=(m_1,\ldots,m_{n-1})\in\mathbb Z_{\geq0}^{n-1}$, with
$\lambda_m=\sum_i m_i\omega_i$. In the usual Euclidean root space
$\{x\in\mathbb R^n:\sum x_i=0\}$ with roots $e_i-e_j$ of squared
length two, set
\[
 \omega_i=\sum_{j=1}^i e_j-\frac in\sum_{j=1}^n e_j,
 \quad \rho_{\rm rt}=\sum_{i=1}^{n-1}\omega_i.
\]
The dimension and the $g_T$-Casimir eigenvalue are
\begin{align}
 d_m&=\prod_{1\leq i<j\leq n}
       \frac{\sum_{q=i}^{j-1}m_q+j-i}{j-i},\label{eq:dim}\\
 C_m&=\frac12(\lambda_m,\lambda_m+2\rho_{\rm rt}).\label{eq:casimir}
\end{align}
These standard representation identities can be read in row coordinates
in \cite{PolychronakosSfetsos}. To spell out the dictionary and constants,
put $\ell_i=\sum_{q=i}^{n-1}m_q$, $\ell_n=0$, and $L=\sum_i\ell_i$.
The highest weight has coordinates $\ell_i-L/n$, and
$2\rho_{{\rm rt},i}=n+1-2i$. Thus \eqref{eq:casimir} says explicitly
\[
 C_m=\frac12\left[\sum_{i=1}^n\ell_i(\ell_i+n+1-2i)-\frac{L^2}{n}\right].
\]
The dimension formula follows from the character alternant
$\det(z_i^{\ell_j+n-j})/\det(z_i^{n-j})$ by letting all $z_i$ tend to
one: the first nonzero alternant terms are the two Vandermonde
determinants in $\ell_j+n-j$ and $n-j$, whose ratio is
\eqref{eq:dim}. The Casimir formula follows by applying the Cartan and
root decomposition of the quadratic Casimir to a highest weight
vector: the Cartan term gives $(\lambda_m,\lambda_m)$, and commuting
each positive-root raising operator past its lowering operator gives
$2(\lambda_m,\rho_{\rm rt})$; the factor $1/2$ is fixed by
\eqref{eq:fundcasimir}. Hence the metric normalization is the one used
in \eqref{eq:H}, rather than an unspecified heat-kernel convention.

Write $k=\sum_i m_i$. Direct scalar multiplication gives
$(\omega_i,\rho_{\rm rt})=i(n-i)/2$. Since
$(\lambda_m,\lambda_m)\geq0$ and $i(n-i)\geq n-1$,
\[
 C_m\geq\frac{n-1}{2}k,\qquad
 d_m\leq(k+1)^{n(n-1)/2}.
\]
The number of $m$ with this sum $k$ is
$\binom{k+n-2}{n-2}\leq(k+1)^{n-2}$: the first $n-2$ entries determine
the last, and each has at most $k+1$ choices. This also covers $n=2$.

Peter--Weyl orthogonality and the Casimir action give the single-edge
heat kernel
\[
 q_s(g)=\sum_m d_m e^{-sC_m}\chi_m(g).
\]
Indeed convolution by $d_m\chi_m$ projects onto the matrix coefficients
of the corresponding irreducible (with the inverse in the kernel
argument if needed), and the heat operator multiplies that subspace by
$e^{-sC_m}$. The series is absolutely uniformly convergent for every
$s>0$: a unitary representation has $|\chi_m(g)|\leq d_m$, and the
above bounds majorize the nonconstant terms by a polynomial in $k$
times $e^{-s(n-1)k/2}$. This also proves finite eigenvalue multiplicities
and divergence of the eigenvalues used in Lemma \ref{lem:ground}.
The constant representation contributes one. For edge $e$ and time $t$,
\begin{align*}
 |q_{\kappa_et}(g)-1|
 &\leq\sum_{k\geq1}(k+1)^D e^{-\kappa_et(n-1)k/2}\\
 &\leq\sum_{k\geq1}\left(2^D e^{-\kappa_et(n-1)/2}\right)^k.
\end{align*}
The second inequality uses $k+1\leq2^k$ for integer $k\geq1$, proved
by induction. At $t=t_*$ the ratio is at most $1/3$, so the sum is at
most $1/2$. Thus $1/2\leq q_{\kappa_et_*}\leq3/2$. The free generator
is the sum of commuting edge generators on the product space, so its
kernel is the product of these $N$ kernels. This proves
\eqref{eq:kernelbounds}.
\end{proof}

\begin{lemma}\label{lem:comparison}
For every continuous $f\geq0$ and $t\geq0$,
\begin{equation}\label{eq:semigroupcomparison}
 e^{-Mt}P_tf\leq e^{-tH}f\leq P_tf.
\end{equation}
Consequently the exact vacuum satisfies
\begin{equation}\label{eq:densitybounds}
 \frac{\max\psi_0}{\min\psi_0}\leq3^Ne^{Mt_*},\qquad
 \frac{\min\rho}{\max\rho}\geq a_*,\qquad
 a_*\leq\rho(U)\leq a_*^{-1}\quad(U\in\mathcal Q).
\end{equation}
\end{lemma}
\begin{proof}
Write $\mathcal A=\sum_e\kappa_e\Delta_e^T$, so
$P_t=e^{t\mathcal A}$. The parabolic maximum principle on this compact
manifold says that a solution of
$(\partial_t-\mathcal A+V)u\geq0$ with nonnegative initial data is
nonnegative when $V\geq0$. One direct proof replaces $u$ by
$e^{-ct}u+\epsilon$ with $c>0$, considers a first negative minimum,
and uses that $\partial_tu\leq0$, $\mathcal Au\geq0$ at that point;
the positive zero-order coefficient makes the differential inequality
impossible. Letting $\epsilon\downarrow0$ proves the result. Smooth
approximation covers continuous initial data.
Apply this principle first to $u=e^{-tH}f$ to get $u\geq0$.
The free solution $v=P_tf$ is also nonnegative by the same argument
with $V=0$. Now
$(\partial_t-\mathcal A+V)v=Vv\geq0$, so $v-u\geq0$.
For $w=e^{-Mt}v$ the corresponding expression is
$(V-M)w\leq0$, so $u-w\geq0$. This proves
\eqref{eq:semigroupcomparison} with no approximation of $V$.

Apply it to $f=\psi_0$ and use
$e^{-tH}\psi_0=e^{-tE_0}\psi_0$. At time $t_*$, letting
$I_0=\int\psi_0d\lambda>0$, Lemma \ref{lem:heat} gives
\[
 e^{(E_0-M)t_*}2^{-N}I_0
 \leq\psi_0(U)\leq e^{E_0t_*}(3/2)^NI_0.
\]
Take the ratio of upper to lower bounds and square to obtain the first
two assertions in \eqref{eq:densitybounds}. Since $\int\rho d\lambda=1$,
$\min\rho\leq1\leq\max\rho$. The ratio bound then gives
$\min\rho\geq a_*\max\rho\geq a_*$ and
$\max\rho\leq\min\rho/a_*\leq a_*^{-1}$, as required.
\end{proof}

\begin{theorem}\label{thm:finite-bounds}
For every simple edge loop,
\begin{equation}\label{eq:momentbounds}
 a_*\leq r_C\leq n^2-a_*(n^2-1),\qquad
 a_*(n^2-1)\leq n^2-r_C\leq n^2-a_*.
\end{equation}
For every such loop carrying a nontrivial center charge,
\begin{equation}\label{eq:quotientbounds}
 \frac{K_C}{2n}\frac{a_*(n^2-1)}{n^2-a_*(n^2-1)}
 \leq Q_C\leq
 \frac{K_C}{2n}\frac{n^2-a_*}{a_*}.
\end{equation}
The full finite-lattice spectral gap, and also its gauge-invariant
restriction above the unique vacuum, obey the genuine lower bound
\begin{equation}\label{eq:gapbound}
 \operatorname{gap}(H)\geq a_*\kappa_*\frac{n^2-1}{2n}>0.
\end{equation}
\end{theorem}
\begin{proof}
The density inequality $\rho\geq a_*$ and Lemma \ref{lem:Haar-loop}
give $r_C\geq a_*$ and
$n^2-r_C\geq a_*(n^2-1)$ by integrating the two nonnegative functions.
These prove all of \eqref{eq:momentbounds}. For a charged loop the
function $r\mapsto(n^2-r)/r$ is strictly decreasing for $r>0$;
substitute the two endpoints into \eqref{eq:chargedquotient} to get
\eqref{eq:quotientbounds}.

For \eqref{eq:gapbound}, first the exact free-product gap is
$g_0=\kappa_*C_F$. To verify the least nonzero single-edge Casimir,
compute
$(\omega_i,\omega_j)=\min(i,j)-ij/n>0$ for $1\leq i,j\leq n-1$.
If $m_i\geq1$, all coefficients in
$C(\lambda_m)-C(\omega_i)$ are nonnegative in the expansion in
fundamental weights. Hence
\[
 C_m\geq C(\omega_i)
 =\frac{i(n-i)(n+1)}{2n}\geq\frac{n^2-1}{2n}.
\]
Equality occurs for a fundamental or antifundamental representation.
On the product space each nonconstant matrix coefficient has at least
one such nonzero edge term; taking a fundamental coefficient on an
edge with $\kappa_e=\kappa_*$ realizes $g_0$. Peter--Weyl expansion
therefore proves the free Poincare inequality
$\mathcal E_\lambda(f,f)\geq g_0\Var_\lambda(f)$.
For any complex $f$, the variational identity
$\Var_\mu(f)=\min_{c\in\mathbb C}\int|f-c|^2d\mu$ gives
\[
 \Var_\mu(f)\leq(\max\rho)\Var_\lambda(f),\qquad
 \mathcal E_\mu(f,f)\geq(\min\rho)\mathcal E_\lambda(f,f).
\]
Consequently
$\mathcal E_\mu(f,f)\geq[(\min\rho)/(\max\rho)]g_0\Var_\mu(f)
\geq a_*g_0\Var_\mu(f)$.
The exact unitary ground-state transform identifies this weighted
Poincare infimum with the spectral gap of $H$. Restricting the infimum
to gauge-invariant functions can only increase it, and the ground
state belongs to that subspace by Lemma \ref{lem:ground}.
\end{proof}

The quotient bounds and the spectral bound have different logical
roles: the upper bound on any one orthogonal trial quotient is an upper
bound on the spectral gap, while a lower bound on that one quotient
alone is not. Equation \eqref{eq:gapbound} is valid because its proof
covers every form-domain vector. All constants in
\eqref{eq:constants} retain volume and coupling dependence. For example,
$a_*$ decreases exponentially in $N$ even with $B$ fixed, and also
decreases with $Bt_*$. Thus none of these inequalities asserts a
strictly positive lower bound after removing a cutoff or taking
infinite volume, and none proves that the limiting gap vanishes.

For a normalization check only, if $V$ is identically zero then
$\psi_0=1$, $r_C=1$, and $m_C=0$ for a simple loop, yielding
$Q_C=K_CC_F$. This is an exact special case of the formula, not a
replacement for the positive-coupling vacuum. The stated interacting
results above apply with all $b_p>0$.

\section{Disjoint parallel plaquettes, covariance Fourier sums, and spectral weight}

Let $p_j$, $j=0,\ldots,m-1$, be simple plaquettes whose physical-edge
supports are pairwise disjoint. Define
\[
 W_j=W_{p_j},\quad f_j=W_j-\mu(W_j),\quad
 c_j(k)=e^{ikj},\quad F_k=\sum_{j=0}^{m-1}c_j(k)f_j.
\]
The mean subtraction is necessary for center-neutral plaquettes. It
gives $\mu(F_k)=0$ for every real $k$ without any translation assumption.
For each edge at most one $f_j$ has a nonzero derivative, so pointwise
\begin{equation}\label{eq:disjointnumerator}
 \Gamma_\kappa(F_k,F_k)
 =\sum_j|c_j(k)|^2\Gamma_\kappa(W_j,W_j)
 =\sum_j\frac{K_j}{2n}(n^2-|W_j|^2),
 \quad K_j=\sum_{e\in p_j}\kappa_e.
\end{equation}
In particular its exact vacuum expectation is independent of $k$:
\begin{equation}\label{eq:Ek}
 E_k:=\mu(\Gamma_\kappa(F_k,F_k))
      =\sum_j\frac{K_j}{2n}(n^2-r_j),\qquad r_j=\mu(|W_j|^2).
\end{equation}
Edge disjointness says nothing about statistical independence in
$\mu$. The full variance is
\begin{equation}\label{eq:Dk}
 D_k:=\mu(|F_k|^2)
 =\sum_{j,l=0}^{m-1}e^{ik(j-l)}C_{jl},\qquad
 C_{jl}=\mu\bigl[(W_j-\mu W_j)\overline{(W_l-\mu W_l)}\bigr].
\end{equation}
All plaquette interactions, including plaquettes meeting more than one
of these supports or connecting them indirectly, remain in $C_{jl}$.
The covariance matrix is Hermitian positive semidefinite because its
quadratic form is the displayed squared $L^2(\mu)$ norm.
In fact $D_k>0$: if $F_k$ were constant, varying an edge unique to $p_j$
would force $W_j$ to be constant, which it is not. Positivity of $\rho$
then gives strictly positive variance. The exact quotient is
\begin{equation}\label{eq:disjointquotient}
 Q_k=\frac{E_k}{D_k}.
\end{equation}

Suppose now the specified finite lattice and its full lists of couplings
are invariant under a translation $\tau$ of order $m$ taking $p_j$ to
$p_{j+1}$ cyclically with their chosen orientations. Such a translation
preserves $H$ and Haar measure and therefore $\psi_0$ and $\mu$ by the
same uniqueness argument as for gauge transformations. There are then
constants $K_j=K$, $r_j=r$ and a cyclic covariance
\[
 C(d)=\mu(f_{l+d}\overline{f_l}),\qquad d\in\mathbb Z/m\mathbb Z,
\]
independent of $l$. For the allowed cyclic wave numbers
$k=2\pi q/m$, $q\in\{0,\ldots,m-1\}$, set
$\widehat C(k)=\sum_{d=0}^{m-1}e^{ikd}C(d)$. Regrouping the full
double sum in \eqref{eq:Dk}, with $j=l+d$, proves
\begin{equation}\label{eq:Fourierquotient}
 D_k=m\widehat C(k),\qquad
 E_k=m\frac{K}{2n}(n^2-r),\qquad
 Q_k=\frac{K(n^2-r)}{2n\widehat C(k)}.
\end{equation}
Here $\widehat C(k)$ is real and strictly positive because it equals
$D_k/m$; $C(-d)=\overline{C(d)}$ also proves reality directly.
There is no algebraic $k^2$ factor in this numerator. Obtaining a
small quotient in this family requires the actual interacting ratio
on the right of \eqref{eq:Fourierquotient} to become small. An assertion
that all covariances vanish, or that the denominator has prescribed
small-$k$ behavior, is not contained in edge disjointness.

For real plaquette energy observables $A_j=n-\Rea W_j$, replace $f_j$
by $A_j-\mu A_j$ and use the same coefficients. All support and covariance
arguments remain exact, but the local numerator is
\begin{equation}\label{eq:real-localnumerator}
 e_j^{\rm real}
 =\frac{K_j}{4}\left[n-\mu\bigl(\Rea\tr(H_{p_j}^2)\bigr)
                       -\frac2n\mu\bigl((\Ima W_j)^2\bigr)\right]
\end{equation}
by \eqref{eq:realenergy}. Thus under the same cyclic invariance
$Q_k=e^{\rm real}/\widehat C_A(k)$ with $\widehat C_A$ the full covariance
transform of $A_j$. Replacing this local numerator by the complex
Wilson numerator would lose terms unless equality is proved in the
particular group and observable.

Finally let $\mathsf P_{H-E_0}(\cdot)$ be the spectral projection measure
of the nonnegative self-adjoint operator $H-E_0$ on the invariant
Hilbert space. The normalized state and its scalar spectral probability
measure are
\begin{equation}\label{eq:spectralmeasure}
 v_k=\frac{\psi_0F_k}{\sqrt{D_k}},\qquad
 \nu_k(A)=\langle v_k,\mathsf P_{H-E_0}(A)v_k\rangle.
\end{equation}
The orthogonality $\langle\psi_0,v_k\rangle=0$ and uniqueness of the
vacuum imply $\nu_k(\{0\})=0$. The spectral theorem and the ground-state
identity prove the exact interacting first moment
\begin{equation}\label{eq:firstmoment}
 \int_{[0,\infty)}\omega\,d\nu_k(\omega)=Q_k=\frac{E_k}{D_k}.
\end{equation}
For every $\epsilon>0$, integrating
$\epsilon\mathbf 1_{[\epsilon,\infty)}(\omega)\leq\omega$ gives
\[
 \nu_k([\epsilon,\infty))\leq Q_k/\epsilon.
\]
This is the exact connection between the computed variation and
low-energy spectral weight. It makes no unproved assertion about
$\widehat C(k)$, a continuum Hamiltonian, an infinite-volume vacuum,
or convergence of the measures $\nu_k$.

\section{Verification and scope}

The companion script \texttt{check\_wilson\_exact.py} uses exact SymPy
arithmetic. On 8 September 2026 it completed with 1,012 exact scalar
assertions, together with exact matrix-entry assertions. It verifies
the basis, Fierz and Casimir constants for $\SU(2)$, $\SU(3)$, $\SU(4)$;
the factor-four metric dictionary; explicit noncommuting $\SU(3)$
positive and negative edge variations and their cross contractions;
real trace energy; a repeated physical edge with opposite signs;
and the one-coordinate product-rule identity underlying
\eqref{eq:groundform}. The proofs above are for all $n\geq2$ and are
not inferred from those finite checks. The script does not compute
$\psi_0$, the interacting covariance matrix, or a continuum spectrum.

All main identities here concern the exact finite interacting
Hamiltonian \eqref{eq:H}. They give explicit trial vectors and their
exact relation to its spectral measure, as well as a positive but
nonuniform finite-volume spectral bound. These proved relations can
be carried into a separate continuum or geometric construction only
through a proved state-space and operator map with controlled limits;
no such map is claimed to be established by the present finite-lattice
calculation.

\begin{thebibliography}{9}
\bibitem{PeterWeyl}
F. Peter and H. Weyl,
\emph{Die Vollst\"andigkeit der primitiven Darstellungen einer
geschlossenen kontinuierlichen Gruppe},
Mathematische Annalen \textbf{97} (1927), 737--755.
\href{https://doi.org/10.1007/BF01447892}{doi:10.1007/BF01447892}.
This is the foundational compact-group completeness theorem used to
identify the complete free spectral resolution.

\bibitem{PolychronakosSfetsos}
A. P. Polychronakos and K. Sfetsos,
\emph{Composing arbitrarily many $SU(N)$ fundamentals},
Nuclear Physics B \textbf{994} (2023), 116314,
\href{https://arxiv.org/abs/2305.19345}{arXiv:2305.19345}.
The row/weight dictionary and dimension/Casimir expressions are given
in Sections 2 and 4, including equation (4.30). Our symbol $n$ is that
paper's $N$; the derivations and all bounds used here retain the
specified anti-Hermitian basis convention.
\end{thebibliography}
\end{document}
